suppose we have an infinitely divisible law on \C^2, supported on the cone {(u,v) \in \C^2; |u|=|v|}. Show that it is actually supported on an additive subgroup contained in the cone

Let
[
K:={(u,v)\in \mathbf C^2:\ |u|=|v|}.
]
We view (\mathbf C^2) as the additive group (\mathbf R^4).

Let (\mu) be infinitely divisible and supported on (K). Since (\mu) is infinitely divisible, it has a convolution square root: there exists a probability measure (\nu) on (\mathbf C^2) such that
[
\mu=\nu*\nu.
]

Set
[
A:=\operatorname{supp}\nu.
]
By the standard support formula for convolution on (\mathbf R^4),
[
\operatorname{supp}(\nu*\nu)=\overline{A+A}.
]
Since (\mu) is supported on (K), and (K) is closed,
[
A+A\subset \overline{A+A}=\operatorname{supp}\mu\subset K.
]

Now define the real quadratic form
[
Q(u,v):=|u|^2-|v|^2,
]
so that (K={Q=0}), and the associated symmetric real bilinear form
[
B\big((u,v),(u',v')\big):=\Re\big(u\overline{u'}-v\overline{v'}\big).
]
Then
[
Q(x+y)=Q(x)+Q(y)+2B(x,y).
]

Take (x\in A). Since (x+x=2x\in A+A\subset K), we get (Q(2x)=0), hence (Q(x)=0). So
[
A\subset K.
]

Next, if (x,y\in A), then (x+y\in A+A\subset K), so
[
0=Q(x+y)=Q(x)+Q(y)+2B(x,y)=2B(x,y),
]
because (Q(x)=Q(y)=0). Therefore
[
B(x,y)=0 \qquad (x,y\in A).
]

So every pair of points of (A) is mutually orthogonal for (B), and every point of (A) is isotropic.

Now let (G) be the additive subgroup generated by (A):
[
G=\Big{\sum_{j=1}^m n_j a_j:\ m\ge1,\ n_j\in\mathbf Z,\ a_j\in A\Big}.
]
For any (g=\sum n_j a_j\in G),
[
Q(g)=\sum_j n_j^2 Q(a_j)+2\sum_{i<j} n_i n_j B(a_i,a_j)=0,
]
since (Q(a_j)=0) and (B(a_i,a_j)=0). Thus
[
G\subset K.
]

Finally, (A+A\subset G), hence
[
\operatorname{supp}\mu=\overline{A+A}\subset \overline{G}.
]
Since (\overline{G}) is again an additive subgroup of (\mathbf C^2), and (K) is closed with (G\subset K), we have
[
\overline{G}\subset K.
]

Therefore (\mu) is supported on the additive subgroup (\overline{G}\subset K).

So any infinitely divisible law on (\mathbf C^2) supported on ({ |u|=|v| }) is in fact supported on an additive subgroup contained in that cone.


what are the additive subgroups of said cone? 

Let
[
K={(u,v)\in \mathbf C^2:\ |u|=|v|}.
]

The answer is:

> **An additive subgroup (G\subset K) is exactly an additive subgroup of a totally isotropic real subspace of (\mathbf C^2).**
> Those real subspaces are only of dimension (0,1,) or (2), and the (2)-dimensional ones are precisely
> [
> P_\lambda^{+}:={(z,\lambda z):z\in \mathbf C},
> \qquad
> P_\lambda^{-}:={(z,\lambda \bar z):z\in \mathbf C},
> \qquad |\lambda|=1.
> ]

So every additive subgroup of the cone is contained in one of those planes (or in an isotropic line inside one of them).

---

## 1. Why every subgroup sits in a totally isotropic real subspace

As before, define
[
Q(u,v)=|u|^2-|v|^2,
]
and the associated real symmetric bilinear form
[
B\big((u,v),(u',v')\big)=\Re!\big(u\overline{u'}-v\overline{v'}\big).
]
Then
[
Q(x+y)=Q(x)+Q(y)+2B(x,y).
]

If (G\subset K) is an additive subgroup, then for any (x,y\in G),
[
x,\ y,\ x+y \in K,
]
so
[
Q(x)=Q(y)=Q(x+y)=0.
]
Hence
[
0=Q(x+y)-Q(x)-Q(y)=2B(x,y),
]
thus
[
B(x,y)=0 \qquad (x,y\in G).
]

Therefore the real span
[
H:=\operatorname{span}_{\mathbf R}(G)
]
satisfies
[
Q(h)=0,\qquad B(h,h')=0\quad \forall h,h'\in H.
]
So (H) is a totally isotropic real subspace for the form (B).

Now (B) has signature ((2,2)) on (\mathbf R^4), so a totally isotropic real subspace has dimension at most (2). Thus (H) has dimension (0,1,) or (2).

---

## 2. Classification of the (2)-dimensional isotropic planes

Let (H\subset \mathbf C^2) be a (2)-dimensional real subspace contained in (K).

Consider the projection (\pi_1:H\to \mathbf C), ((u,v)\mapsto u).
If ((0,v)\in H), then (|v|=|0|=0), so (v=0). Hence (\pi_1) is injective. Since both (H) and (\mathbf C\cong \mathbf R^2) are 2-dimensional over (\mathbf R), (\pi_1) is an isomorphism of real vector spaces.

So (H) is the graph of a real-linear map (T:\mathbf C\to \mathbf C):
[
H={(z,Tz):z\in \mathbf C}.
]
Because (H\subset K), for all (z),
[
|Tz|=|z|.
]
So (T) is a real-linear isometry of (\mathbf C\simeq \mathbf R^2), i.e. (T\in O(2)).

Every real-linear orthogonal map on (\mathbf C) is either
[
Tz=\lambda z \quad\text{or}\quad Tz=\lambda \bar z
\qquad (|\lambda|=1).
]
Therefore
[
H=P_\lambda^{+}={(z,\lambda z):z\in\mathbf C}
]
or
[
H=P_\lambda^{-}={(z,\lambda\bar z):z\in\mathbf C}.
]

These are indeed contained in (K), since
[
|\lambda z|=|z|,\qquad |\lambda \bar z|=|z|.
]

---

## 3. The 1-dimensional isotropic subspaces

A real line (\mathbf R\cdot (u,v)) is contained in (K) iff (|u|=|v|).
So the 1-dimensional cases are exactly the real lines generated by nonzero points of the cone.

Equivalently, every such line is contained in some (P_\lambda^+) or (P_\lambda^-).

---

## 4. Final classification of additive subgroups

Putting it together:

### Every additive subgroup (G\subset K) is contained in one of:

* ({0}),
* a real isotropic line (\mathbf R\cdot (u,v)) with (|u|=|v|),
* a plane (P_\lambda^+) or (P_\lambda^-), where (|\lambda|=1).

Conversely, **any additive subgroup** of any one of these subspaces is contained in (K), because the whole subspace is contained in (K).

So the clean classification is:

[
\boxed{
G\subset K \text{ additive }
\iff
\exists,|\lambda|=1\ \text{such that } G\subset P_\lambda^+
\text{ or } G\subset P_\lambda^-.
}
]

where
[
P_\lambda^+={(z,\lambda z):z\in\mathbf C},
\qquad
P_\lambda^-={(z,\lambda\bar z):z\in\mathbf C}.
]

Inside one such plane, (G) can be any additive subgroup at all: dense, discrete, a lattice, a line, etc.

---

## 5. Examples

[
{(z,z):z\in\mathbf Z[i]}\subset P_1^+
]
is an additive subgroup of the cone.

[
{(z,\bar z):z\in\mathbf C}=P_1^-
]
is itself an additive subgroup of the cone.

[
\mathbf R\cdot(1,e^{i\theta})
]
is a 1-dimensional additive subgroup of the cone.

---

If you want, I can also phrase this in a more representation-theoretic way, as “the cone contains no nontrivial additive structure except along graphs of orthogonal real-linear maps (\mathbf C\to\mathbf C).”

